% !TEX program = xelatex
\documentclass{misisarticle}

\grnti{20.23.17, 28.29.07, 28.17.19}
\asjc{1710, 1708, 2611}
\oecd{1.2, 1.1}

\title{Сравнительный анализ и классификация моделей балансировки нагрузки в распределённых СУБД}

% TODO: аффилиация, город и email.
\misisauthor{Хлопцев Артём Евгеньевич}{}{}{}

% TODO: аннотация.
\abstract{}
\keywords{балансировка нагрузки; распределённые СУБД; NewSQL; шардирование; консенсус Raft; горячие сегменты.}

\begin{document}
\maketitle

\section{Введение}
% TODO: актуальность, проблема, цель и границы исследования.

\section{Анализ работ}
% TODO: анализ литературы, классификация подходов и исследовательский пробел.

\section{Постановка задачи}
% TODO: переменные размещения, модель нагрузки, целевые функции и ограничения.

\section{Описание методов}
% TODO: базовые стратегии, предлагаемая модель и протокол проверки.

\begin{thebibliography}{9}

\bibitem{curino2010}
Curino C., Jones E. P. C., Zhang Y., Madden S.
Schism: A Workload-Driven Approach to Database Replication and Partitioning //
PVLDB. 2010. Vol.~3, No.~1--2. P.~48--57.
\url{https://doi.org/10.14778/1920841.1920853}.

\bibitem{lee2021}
Lee S., Guo Z., Sunercan O. et al.
Shard Manager: A Generic Shard Management Framework for Geo-distributed Applications //
SOSP 2021. P.~553--569.
\url{https://doi.org/10.1145/3477132.3483546}.

\bibitem{taft2020}
Taft R., Sharif I., Matei A. et al.
CockroachDB: The Resilient Geo-Distributed SQL Database //
SIGMOD 2020. P.~1493--1509.
\url{https://doi.org/10.1145/3318464.3386134}.

\bibitem{antonov2023}
Антонов И.~А., Исаченко М.~К.
Методы машинного обучения для обнаружения вторжений в сетях интернета вещей //
Искусственный интеллект в промышленных, коммерческих, медицинских и финансовых приложениях:
сборник статей научно-технического семинара студентов.
Вып.~1. М.: НИТУ МИСИС, 2023. С.~9--15.
\url{https://sadekov.su/Articles/kik_2023.pdf#page=11}.

\bibitem{sosedka2024}
Соседка А.~В., Ибрагимов П.~И.
Количественный и качественный анализ аудитории Telegram в разрезе рекомендаций с использованием больших языковых моделей //
Искусственный интеллект в промышленных, коммерческих, медицинских и финансовых приложениях.
М.: НИТУ МИСИС, 2024. С.~116--120.
\url{https://sadekov.su/Articles/kik_2024_1.pdf#page=117}.

\bibitem{tarkhanov2020}
Tarkhanov I.~A., Akimova G.~P., Pashkin M.~A., Solovyev A.~V.
Modelling the Methodology to Assess the Effectiveness of Distributed Information Systems //
ASTESJ. 2020. Vol.~5, No.~1. P.~86--92.
\url{https://astesj.com/v05/i01/p11}.

\end{thebibliography}
\end{document}
